%----------------------------------------------------------------------------------------
%	PACKAGES AND OTHER DOCUMENT CONFIGURATIONS
%----------------------------------------------------------------------------------------

\documentclass[letterpaper]{resume} % Use the custom resume.cls style
\usepackage{nopageno}
\usepackage[left=0.40in,top=0.3in,right=0.75in,bottom=0.3in]{geometry} % Document margins
\usepackage{fontawesome}
\usepackage{times}
    \newcommand{\tab}[1]{\hspace{.2667\textwidth}\rlap{#1}}
    \newcommand{\itab}[1]{\hspace{0em}\rlap{#1}}
%\usepackage[german]{babel}
\usepackage{ulem}
\usepackage{enumitem}
    \setitemize{noitemsep,topsep=0pt,parsep=0pt,partopsep=0pt}
    
\name{José Alfredo Rosas Córdova} % Your name 
\address{\faEnvelope{ jarosascordova@gmail.com}\hspace{0.75em}\faPhone{ 961 279 9900}}
\address{\faMapMarker{Alondra 115, Col. Cuauhtémoc, 66450 San Nicolás de los Garza}} % Your address
\address{\faBirthdayCake{ 22/10/1999}\hspace{0.75em}\faGlobe{ Guadalajara, México}}

\begin{document}
    \begin{rSection}{Formación Académica}
        {\bf Ingeniería en Aeronáutica} \hfill {\itshape 08/2017 --- 06/2023} 
        \\{ \textit{Universidad Autónoma de Nuevo León}}
        Promedio de 96.5
    
        {\bf Estancias Combinadas de Estudios y Prácticas (DAAD-KOSPIE)} \hfill {\itshape 09/2021 --- 09/2022} 
        \\{\bf Ingeniería Mecánica}
        \\ { \textit {Technische Universität Braunschweig}}

        {\bf Bachillerato General} \hfill {\itshape 08/2014 --- 07/2017} 
        \\{ \textit{Colegio de Bachilleres de Chiapas - Plantel 33 Polyforum}}
    \end{rSection}
    
    \begin{rSection}{Idiomas}
        \begin{tabular}{ @{} >{\bfseries}l @{\hspace{6ex}} l }
            Español & Lengua materna\\
            Inglés & TOEFL ITP -- 677 Puntos (C1)\\
            Alemán & TestDaF -- TDN 5 (C1)
        \end{tabular}
    \end{rSection}
    
    \begin{rSection}{Habilidades}
        \begin{tabular}{ @{} >{\bfseries}l @{\hspace{6ex}} l }
            CAD \& CAM: & SOLIDOWRKS, CATIA V5, OnShape, Autodesk Fusion360\\
            Programación: & Python -- Competente\\ & MATLAB -- Competente\\ & C++ -- Fundamentos\\
                Modelado \& Simulación: & ANSYS Workbench, Fluent, Simulink, Xfoil, OpenVSP\\
            Otros: & LaTeX, MS Office, GitHub, Jira\\
        \end{tabular}
    \end{rSection}
    
    \begin{rSection}{Experiencia Laboral}
        {\bf Ingeniero de Mecánico} \hfill{\itshape 01/2023 --- 04/2023} 
        \\{\textit{Vertiv México -- Monterrey, México}}
        \begin{itemize}
            \item Diseño mecánico de tuberías y sheet metal.
            \item Procesamiento de datos de pruebas de laboratorio.
            \item Asistencia en construcción de prototipos.
        \end{itemize}
        
        {\bf Ingeniero de Proyecto - Desarrollo de Aeronaves no Tripuladas} \hfill{\itshape 01/2023 --- 04/2023} 
        \\{\textit{Spacelab MX -- Monterrey, México}}
        \begin{itemize}
            \item Evaluación de la viabilidad de conceptos de misión, selección y posterior desarollo del mismo.
            \item Dimensionamiento inicial de la aeronave y estudios preliminares de desempeño de vuelo.
            \item Diseño estructural preliminar de la aeronave para la fase de prototipado del proyecto.
            \item Pruebas de integración, en tierra y en vuelo de la aeronave del banco de pruebas
        \end{itemize}
    
        {\bf Prácticante de Diseño Mecánico y Pruebas de Sistema, Departamento de Mecatrónica} \hfill{\itshape 04/2022 --- 08/2022} 
        \\{\textit{Bosch Rexroth AG -- Stuttgart, Alemania}}
        \begin{itemize}
            \item Diseño de elementos internos de hardware y tren de potencia del vehículos autónomos no tripulados.
            \item Prototipado de soluciones de diseño propuestas acompañadas de un proceso de la validación en un contexto de desarrollo.
            \item Colaboración con otros equipos del departamento para la evaluación de los requisitos de diseño y el desarrollo de funciones en tanto hardware y software.
            \item Evaluación, análisis y detección de fallos del desempeño de conducción de los UGV, así como de la capa de comunicación entre los vehículos y el servicio de gestión de flotas basándose en datos obtenidos de pruebas al sistema.
        \end{itemize}
    \end{rSection}
    
    
    %--------------------------------------------------------------------------------
    %    Projects And Seminars
    %-----------------------------------------------------------------------------------------------
    \clearpage
    \begin{rSection}{Investigación}
        {\bf Becario, Laboratorio de Aerodinámica} \hfill {\itshape desde 07/2020} 
        \\{\textit{Centro de Investigación e Innovación en Ingeniería Aeronáutica}}
        \begin{itemize}
            \item Programa de Financiamiento a la Investigación Científica y Tecnológica de la UANL \hfill {\itshape 05/2021}
            \begin{itemize}
                \item Estudio de Aeronaves no tripuladas. Número de proyecto: IT1814-21
            \end{itemize}

            \item XXII Verano de Investigación Científica y Tecnológica de la UANL\\ PROVERICYT \hfill {\itshape 07/2020}
            \begin{itemize}
                \item Nuevos métodos de control para aeronaves pequeñas
            \end{itemize}
            \item Servicio social \hfill {\itshape since 09/2022}
            \begin{itemize}
                \item Mantenimiento de equipos de fabricación por control numérico computarizado (CNC).
                \item Desarrollo de aeronaves autónomas.
            \end{itemize}
        \end{itemize}

    {\bf Prestador de Servicio Social, Laboratorio de Diseño de Estructuras Aeroespaciales} \hfill {\itshape 09/2022 --- 12/2022} 
        \\{\textit{Centro de Investigación e Innovación en Ingeniería Aeronáutica}}
            \begin{itemize}
                \item Mantenimiento a equipos de manufactura del laboratorio. Primordialmente enfocado a equipos por control numérico (CNC).
                \item Desarrollo y fabricación de aeronaves autónomas con un enfoque a aplicaciones académicas.
            \end{itemize}
    \end{rSection}
    
    
    %----------------------------------------------------------------------------------------
    %	WORK EXPERIENCE SECTION
    %----------------------------------------------------------------------------------------
    % \begin{rSection}{Work Experience}
    %     {\bf Work Station x,Strange Place x} \hfill {\em} 
    %     \\{\textit{ Engineer}}
    %     \\- xxx
    %     \\- xxx
    % \end{rSection}
    
    
    %----------------------------------------------------------------------------------------
    %	Extracurricular Section
    
    %----------------------------------------------------------------------------------------
    \begin{rSection}{Actividades Extracurriculares}
        {\bf Ingeniero en Jefe, Axios Aero Design} \hfill {\itshape 07/2021 --- 06/2022} 
        \\{\textit{Capítulo Estudiantil de SAE Aero Design}}
        \begin{itemize}
            \item Estudios conceptuales y de diseño preliminar de la aeronave. Abarcando áreas de aerodinámica, estructuras, propulsión eléctrica, y estabilidad y control de vuelo. \item Planificación y ejecución de proyectos basados en metodologías de ingeniería de sistemas.
            \item  Asistencia en el diseño detallado de componentes estructurales de aeronaves como ala, fuselaje y empenaje.
        \end{itemize}

        {\bf Encargado del Reporte Técnico, Axios Aero Design} \hfill {\itshape 09/2020 --- 06/2021} 
        \\{\textit{Capítulo Estudiantil de SAE Aero Design}}
        \begin{itemize}
            \item Producción del reporte técnico del proyecto, consistente de la documentación del proceso de diseño de ingeniería, así como dibujos técnicos que describan el producto final.
        \end{itemize}
    \end{rSection}
    
    \begin{rSection}{Premios}
        {\bf SAE Aero Design México 2022} \hfill {\itshape 03/2022}
        \\{\textit{2. \textordmasculine~Lugar General}}
        \\{\textit{3. \textordmasculine~Lugar en Reporte Técnico}}
    
        {\bf MathWorks Minidrone Competition Mexico 2021} \hfill {\itshape 08/2021}
        \\{\textit{1. \textordmasculine~Lugar}}
        
        {\bf SAE Aero Design México 2020} \hfill {\itshape 10/2020}
        \\{\textit{6. \textordmasculine~Lugar General}}
        \\{\textit{4. \textordmasculine~Lugar en Reporte Técnico}}
    \end{rSection}

    \begin{rSection}{}
        \vspace{7.5em}
        \begin{minipage}{0.5\linewidth}
            \uline{San Nicolás de los Garza, NL \today}\\
            \textit{Lugar, Fecha}
        \end{minipage}
        \begin{minipage}{0.5\linewidth}
            \uline{José Alfredo Rosas Córdova\hspace{5em}}\\
            \textit{Firma}
        \end{minipage}
    \end{rSection}
\end{document}